\chapter{Human-Review Protocol and Supporting Materials}
\label{app:human-review}

\section{Review Sample and Presentation}

The retained review package contains 64 clips, balanced across the four action classes
and spread across recordings. It is an audit sample, not a population-random estimate of
annotation accuracy. Fixed blind IDs and sheet order support correspondence between the
material presented and returned ratings. The browser interface supplied compatible WebM
playback copies while retaining the original MP4 files. An additional encoding can
soften visual detail, so the originals remain relevant when investigating presentation
artefacts. Package checks addressed file hashes, decoding, frame counts, cadence and
interface logic; they did not themselves establish annotation correctness.\ev{Interface
preparation and preserved inputs}{\EvidenceRef{E048}{}}

One reviewer returned complete ratings for all 64 clips. The submitted rating file was
checked against the blind identifiers, required fields, allowed responses and reasons
for flagged cases. The first assessment was complete at the reporting date. The flagged
cases awaited adjudication.

\section{Single-Reviewer Results}

\begin{table}[htbp]
\caption{Reviewer Agreement With Recorded Outcomes for 64 Clips}\label{tab:human-review-outcomes}
\centering

\begin{tabular}{lrrrr}
\toprule Action & Agree & Disagree & Indeterminate & Total \\
\midrule
Pick &13&0&3&16\\
Insert &14&1&1&16\\
Remove &15&0&1&16\\
Place &15&1&0&16\\
\midrule
Total &57&2&5&64\\
\bottomrule
\end{tabular}
\end{table}

The reviewer agreed with the recorded outcome on 57 of 59 determinate clips (96.6\%) and
disagreed on two (3.4\%). Five of the 64 clips had indeterminate outcomes (7.8\%).
Across the full sample, agreement was 57 of 64 (89.1\%). The 96.6\% estimate is
conditional on a judgeable outcome. The 89.1\% estimate includes all sampled clips.
Equal sampling of each action class limits generalisation of either percentage to the
full dataset.\ev{Completed ratings, counts and limitations}{\EvidenceRef{E049}{}}

\section{Boundary and Presentation Flags}

Boundary ratings were 60 acceptable, three questionable and one indeterminate.
Visibility was rated adequate for all 64 clips. Four clips received possible recording
or presentation artefact flags. Because flag categories overlap, their sum is not the
number of affected clips: ten unique clips require follow-up. The flag list retains
reasons and exact blind IDs. An adequate visibility rating also does not imply that the
observed interval unambiguously establishes the action outcome; five outcomes remained
indeterminate despite that visibility assessment.

Adjudication remained incomplete. The original clips and annotation boundaries would
help distinguish an ambiguous task outcome from a boundary or playback problem. The
evaluation used the original labels for all cases, including those flagged. No
thresholds or predictions were changed following the review.

\section{Interpretation and Reviewer Independence}

One reviewer supplied all ratings, so inter-rater agreement could not be estimated. The
results describe agreement between that reviewer and the recorded labels for the sampled
clips. Establishing inter-rater reliability would require a second reviewer who had not
seen the first reviewer's ratings or interpretations.

I had access to generic AI-generated wording examples, but these were not additional
clip-level ratings. The reviewer supplied every rating used in the analysis. Checks of
the submitted fields and totals supplied no independent judgement about the clips.
